Para aplicar el método de Lemke se necesita una solución básica que cumpla el criterio de optimalidad, lo que exige un problema de máximos sin restricciones de igualdad. Se desdobla la igualdad en dos desigualdades y se reformula como maximización en forma estándar:

\begin{equation}
\begin{split}
\mbox{max. } z = -10x_1 - 20x_2 - 30x_3\\
s.a.:\\
4x_1 + 2x_2 + 8x_3 + h_1 = 1000\\
x_2 + x_3 + h_2 = 500\\
x_2 + x_3 - h_3 = 500\\
x_1,\,x_2,\,x_3,\,h_1,\,h_2,\,h_3\geq 0
\end{split}
\end{equation}

Las variables de holgura forman una base que cumple el criterio de optimalidad (todos los costes reducidos son no positivos) aunque no es factible ($h_3 < 0$), que es el punto de partida del método:

\begin{center}
\begin{tabular}{c|c|cccccc|}
  & $s$ & $x_1$ & $x_2$ & $x_3$ & $h_1$ & $h_2$ & $h_3$\\
 & $0$ & $-10$ & $-20$ & $-30$ & $0$ & $0$ & $0$\\
\hline
$h_1$ & $1000$ & $4$ & $2$ & $8$ & $1$ & $0$ & $0$\\
$h_2$ & $500$ & $0$ & $1$ & $1$ & $0$ & $1$ & $0$\\
$h_3$ & $-500$ & $0$ & $-1$ & $-1$ & $0$ & $0$ & $1$\\
\hline
 & $10000$ & $-10$ & $0$ & $-10$ & $0$ & $0$ & $-20$\\
\hline
$h_1$ & $0$ & $4$ & $0$ & $6$ & $1$ & $0$ & $2$\\
$h_2$ & $0$ & $0$ & $0$ & $0$ & $0$ & $1$ & $1$\\
$x_2$ & $500$ & $0$ & $1$ & $1$ & $0$ & $0$ & $-1$\\
\hline
\end{tabular}
\end{center}

En la solución óptima del problema original:

\begin{itemize}
    \item $z = -s = 10000$
    \item $x_1 = 0$
    \item $x_2 = 500$
    \item $x_3 = 0$
\end{itemize}
